
前面的 $q\bar q\to gg$ 已经说明：一旦外态含胶子，非阿贝尔三胶子顶角不再是可有可无的修正，而是 Ward 恒等式所要求的树级组成部分。要把强子对撞机中最常见的 Born 硬核补齐，还需要再看两个含初态胶子的通道 $qg\to qg$ 与 $gg\to gg$。它们的意义不只是“再多算两个过程”：前者同时包含基本表示与伴随表示的颜色代数，后者则完全由胶子自相互作用组成，因而把 $C_F$ 与 $C_A$ 在可观测散射核中的不同角色直接显露出来。

先保留夸克质量 $m_q$，取
\begin{equation*}
 q(p_1,i)+g(k_1,a)\to q(p_2,j)+g(k_2,b),
 \qquad
 p_1^2=p_2^2=m_q^2c^2,
 \qquad k_1^2=k_2^2=0,
\end{equation*}
并定义
\begin{equation*}
 s=(p_1+k_1)^2,
 \qquad
 t=(p_1-p_2)^2,
 \qquad
 u=(p_1-k_2)^2,
 \qquad s+t+u=2m_q^2c^2.
\end{equation*}
树级振幅由 $s$、$u$ 道夸克交换以及 $t$ 道胶子交换三部分组成。把外胶子偏振矢量记为 $\varepsilon_\mu(k_1)$ 与 $\varepsilon_\nu^*(k_2)$，则可写成
\begin{equation}
\boxed{
\begin{aligned}
\mathcal M_{qg}(m_q)
=g_s^2\,\bar u(p_2)\Bigg[&
 T^bT^a\,\gamma^\nu
 \frac{\slashed p_1+\slashed k_1+m_qc}{s-m_q^2c^2}\gamma^\mu
 +T^aT^b\,\gamma^\mu
 \frac{\slashed p_1-\slashed k_2+m_qc}{u-m_q^2c^2}\gamma^\nu\\
&+\frac{\ii f^{abc}T^c}{t}\,
 \Gamma^{\rho\mu\nu}(p_1-p_2,k_1,-k_2)\gamma_\rho
\Bigg]u(p_1)\,
\varepsilon_\mu(k_1)\varepsilon_\nu^*(k_2).
\end{aligned}}
\label{eq:qcd-qg-amplitude-dp94}
\end{equation}
其中 $\Gamma^{\rho\mu\nu}$ 是前文已经由 $F^a_{\mu\nu}F^{a\mu\nu}$ 推出的三胶子顶角张量。第一条腿属于 $t$ 道内部胶子；因为三胶子规则把三条动量都定义为流入顶角，所以它携带的是 $p_1-p_2=k_2-k_1$，并满足 $(p_1-p_2)+k_1-k_2=0$。这条动量和是使用顶角公式前必须先做的检查。三项的相对系数由同一个协变导数与场强张量固定；若去掉第三项，纵向偏振替换不会相消。

有限 $m_q$ 的振幅由\eqnrefs{eq:qcd-qg-amplitude-dp94} 完整定义。强子对撞机的轻 parton Born 核进一步取严格连续极限 $m_q\to0$；此时 $s+t+u=0$。对初态夸克的 $N_c$ 个颜色和两个自旋、初态胶子的 $N_c^2-1$ 个颜色和两个物理偏振平均，并对末态求和，得到无质量部分子极限的闭式
\begin{equation}
\boxed{
\overline{|\mathcal M(qg\to qg)|^2}
=g_s^4(s^2+u^2)
\left[
 \frac1{t^2}-\frac{C_F}{C_A}\frac1{su}
\right],
\qquad
C_F=\frac{N_c^2-1}{2N_c},\quad C_A=N_c.}
\label{eq:qcd-qg-M2-dp94}
\end{equation}
这里第一项是前向 $t$ 道胶子交换产生的强增强，第二项来自夸克交换与非阿贝尔颜色排序之间的相干组合。对 QCD 的 $N_c=3$，$C_F/C_A=4/9$，恢复常用的
\begin{equation*}
 \overline{|\mathcal M|^2}
 =g_s^4(s^2+u^2)
 \left(\frac1{t^2}-\frac{4}{9su}\right).
\end{equation*}
由于物理区 $s>0$、$t<0$、$u<0$，第二项实际上为正；写成 $-1/(su)$ 的形式是为了保持交叉对称的 Mandelstam 结构。

纯胶子散射
\begin{equation*}
 g(k_1,a)+g(k_2,b)\to g(k_3,c)+g(k_4,d)
\end{equation*}
包含 $s,t,u$ 三个三胶子交换图以及一个四胶子接触图。用 Jacobi 恒等式把接触项分配到三个交换道后，完整树级振幅可按伴随表示颜色张量写成
\begin{equation}
\boxed{
 \mathcal M^{abcd}
 =g_s^2\left[
 f^{abe}f^{cde}\mathcal K_s
 +f^{ace}f^{bde}\mathcal K_t
 +f^{ade}f^{bce}\mathcal K_u
 \right],}
\label{eq:qcd-gg-color-amplitude-dp94}
\end{equation}
其中 $\mathcal K_s,\mathcal K_t,\mathcal K_u$ 包含相应交换图与四胶子接触项的洛伦兹结构。单独任何一项都依赖规范选择；只有它们的完整和在外腿上满足 Ward 恒等式。对一般 $SU(N_c)$ 完成颜色、偏振求和与初态平均后得到
\begin{equation}
\boxed{
\overline{|\mathcal M(gg\to gg)|^2}
=\frac{4g_s^4C_A^2}{N_c^2-1}
\left(
3-\frac{su}{t^2}-\frac{st}{u^2}-\frac{tu}{s^2}
\right).}
\label{eq:qcd-gg-M2-dp94}
\end{equation}
对 $N_c=3$，整体颜色系数成为 $9/2$。与 $qq'$ 或 $qg$ 相比，\eqnrefs{eq:qcd-gg-M2-dp94} 同时在 $t\to0$ 与 $u\to0$ 出现交换增强，反映了两个完全相同的无质量规范玻色子既可沿前向也可沿后向交换。式中的 $C_A^2$ 还直接表明：胶子自相互作用的强度由伴随表示 Casimir 控制，而不是由夸克的 $C_F$ 控制。

对于任意无质量 $2\to2$ Born 过程，统一相空间关系为
\begin{equation}
 \frac{\dd\sigma}{\dd t}
 =\frac{\hbar^2}{16\pi s^2}\overline{|\mathcal M|^2}.
 \label{eq:qcd-partonic-dsdt-master-dp94}
\end{equation}
因此 \eqnrefs{eq:qcd-qg-M2-dp94,eq:qcd-gg-M2-dp94} 已经足以直接生成角分布。\figref{fig:qcd-qg-gg-angular-dp94} 把 $N_c=3$ 时的无量纲 Born 核除以各自在 $90^\circ$ 的值；这样既去掉共同的 $g_s^4$ 与 $1/s$ 尺度，又保留由交换道和颜色结构决定的角形状。

\begin{figure}[htbp]
 \centering
 \includegraphics[width=.74\textwidth]{qcd_qg_gg_angular_dp94.pdf}
 \caption{QCD 中 $qg\to qg$ 与 $gg\to gg$ 的树级无量纲角分布核，均按 $90^\circ$ 归一化。$qg$ 通道主要显示 $t$ 道前向增强；$gg$ 通道还因相同外胶子而出现明显的后向交换增强。曲线由\eqnrefs{eq:qcd-qg-M2-dp94,eq:qcd-gg-M2-dp94} 逐点计算。}
 \label{fig:qcd-qg-gg-angular-dp94}
\end{figure}

\begin{derivation}[从\eqnrefs{eq:qcd-qg-amplitude-dp94}到\eqnrefs{eq:qcd-qg-M2-dp94}]
先取轻夸克极限 $m_q\to0$。这一极限下三张图的相对系数仍由规范不变作用量固定。单张图本身具有规范依赖，因此偏振求和先保留两种横向物理偏振；只有完整振幅相加并满足 Ward 恒等式以后，才可以使用相应的协变化简。

先固定质心系，使每个四动量和物理偏振都能逐分量检查。
令 $Q\equiv\sqrt{s}/2$，并取
\begin{align*}
 p_1^\mu&=Q(1,0,0,1),&
 k_1^\mu&=Q(1,0,0,-1),\\
 p_2^\mu&=Q(1,\sin\theta,0,\cos\theta),&
 k_2^\mu&=Q(1,-\sin\theta,0,-\cos\theta).
\end{align*}
于是
\begin{equation*}
 s=4Q^2,
 \qquad
 t=-2Q^2(1-\cos\theta),
 \qquad
 u=-2Q^2(1+\cos\theta),
 \qquad s+t+u=0.
\end{equation*}
记 $c_\theta\equiv\cos\theta$。入射胶子沿 $-z$ 方向，可选两条实的横向偏振
\begin{equation*}
 \varepsilon_{1x}^\mu=(0,1,0,0),
 \qquad
 \varepsilon_{1y}^\mu=(0,0,1,0),
\end{equation*}
而出射胶子沿 $(-\sin\theta,0,-\cos\theta)$，可选
\begin{equation*}
 \varepsilon_{2\parallel}^\mu=(0,-c_\theta,0,\sin\theta),
 \qquad
 \varepsilon_{2y}^\mu=(0,0,1,0).
\end{equation*}
逐项有 $k_i\cdot\varepsilon_{i\lambda}=0$ 与 $\varepsilon_{i\lambda}^2=-1$。因此这四种偏振组合恰好穷尽两条外胶子的物理态，不含时间样或纵向态。

三张图的 Dirac 核分别定义为
\begin{align*}
 \mathscr A_s^{\mu\nu}
 &=\gamma^\nu\frac{\slashed p_1+\slashed k_1}{s}\gamma^\mu,\\
 \mathscr A_u^{\mu\nu}
 &=\gamma^\mu\frac{\slashed p_1-\slashed k_2}{u}\gamma^\nu,\\
 \mathscr A_t^{\mu\nu}
 &=\frac1t\Gamma^{\rho\mu\nu}(p_1-p_2,k_1,-k_2)\gamma_\rho.
\end{align*}
这里第一条三胶子腿的动量必须是 $p_1-p_2=k_2-k_1$，因为前文三胶子顶角约定三条动量全部流入顶角：
\begin{equation*}
 (p_1-p_2)+k_1-k_2=0.
\end{equation*}
这也是正文\eqnrefs{eq:qcd-qg-amplitude-dp94} 中动量方向的直接检查。

接着计算六个独立的 Dirac--偏振收缩。
对 $x,y\in\{s,u,t\}$ 定义
\begin{equation*}
 K_{xy}(\theta)
 \equiv
 \sum_{r_1,r_2}
 \sum_{\lambda_1,\lambda_2}
 \Bigl[
 \bar u_{r_2}(p_2)\mathscr A_x^{\mu\nu}u_{r_1}(p_1)
 \varepsilon_{1\lambda_1,\mu}
 \varepsilon_{2\lambda_2,\nu}
 \Bigr]
 \Bigl[
 \bar u_{r_2}(p_2)\mathscr A_y^{\alpha\beta}u_{r_1}(p_1)
 \varepsilon_{1\lambda_1,\alpha}
 \varepsilon_{2\lambda_2,\beta}
 \Bigr]^* .
\end{equation*}
利用无质量自旋和
\begin{equation*}
 \sum_r u_r(p)\bar u_r(p)=\slashed p
\end{equation*}
后，每个 $K_{xy}$ 都只是 $4\times4$ Dirac 矩阵的有限迹。含三胶子顶角的项还可以先把 Lorentz 三指标降成一个普通 Dirac 矩阵。令 $q=p_1-p_2=k_2-k_1$，则
\begin{align*}
 &\varepsilon_1^\mu\varepsilon_2^\nu
 \Gamma^{\rho}{}_{\mu\nu}(q,k_1,-k_2)\\
 &\qquad={}
 -2(k_1\!\cdot\varepsilon_2)\varepsilon_1^\rho
 +(\varepsilon_1\!\cdot\varepsilon_2)(k_1+k_2)^\rho
 -2(k_2\!\cdot\varepsilon_1)\varepsilon_2^\rho.
\end{align*}
这一步直接来自三胶子顶角的三项，并只用了
$k_1\cdot\varepsilon_1=k_2\cdot\varepsilon_2=0$ 与 $q=k_2-k_1$。对上面选定的四种物理偏振，它进一步化为
\begin{align*}
 \mathscr A_t^{(x,y)}
 &=-\frac{2Q\sin\theta}{t}\slashed\varepsilon_y,
 &
 \mathscr A_t^{(y,\parallel)}
 &=-\frac{2Q\sin\theta}{t}\slashed\varepsilon_y,\\
 \mathscr A_t^{(y,y)}
 &=-\frac{\slashed k_1+\slashed k_2}{t},
 &
 \mathscr A_t^{(x,\parallel)}
 &=\frac{-2Q\sin\theta\slashed\varepsilon_{1x}
 +c_\theta(\slashed k_1+\slashed k_2)
 -2Q\sin\theta\slashed\varepsilon_{2\parallel}}{t}.
\end{align*}
因此 $t$ 道也已被化成普通的有限伽马矩阵链，不再剩下未计算的三胶子张量。作为第一项示范，$s$ 道平方变成
\begin{align*}
 K_{ss}
 ={}&\sum_{\lambda_1,\lambda_2}
 \Tr\!\left[
 \slashed p_2
 \slashed\varepsilon_2
 \frac{\slashed p_1+\slashed k_1}{s}
 \slashed\varepsilon_1
 \slashed p_1
 \slashed\varepsilon_1
 \frac{\slashed p_1+\slashed k_1}{s}
 \slashed\varepsilon_2
 \right].
\end{align*}
把 $p_1+k_1=(2Q,0,0,0)$ 与上面的四条物理偏振逐一代入，并反复使用
\begin{equation*}
 \{\gamma^\mu,\gamma^\nu\}=2\eta^{\mu\nu}I_4,
 \qquad
 \Tr(\slashed a\slashed b\slashed c\slashed d)
 =4\bigl(a\!\cdot b\,c\!\cdot d-a\!\cdot c\,b\!\cdot d+a\!\cdot d\,b\!\cdot c\bigr),
\end{equation*}
四种偏振组合相加给
\begin{equation*}
 K_{ss}=4(1+c_\theta).
\end{equation*}

再显式算一个含三胶子顶角的干涉项 $K_{st}$。定义 Dirac 共轭矩阵
\begin{equation*}
 \overline{\mathscr A}_t
 \equiv\gamma^0\mathscr A_t^\dagger\gamma^0,
\end{equation*}
则固定一组物理偏振 $(\lambda_1,\lambda_2)$ 时的自旋和为
\begin{equation*}
 K_{st}^{(\lambda_1,\lambda_2)}
 =\Tr\!\left[
 \slashed p_2\,
 \mathscr A_s^{(\lambda_1,\lambda_2)}
 \slashed p_1\,
 \overline{\mathscr A}_t^{(\lambda_1,\lambda_2)}
 \right].
\end{equation*}
把上一段已经化简出的四个 $\mathscr A_t$ 与
$\mathscr A_s=\slashed\varepsilon_2(\slashed p_1+\slashed k_1)\slashed\varepsilon_1/s$
逐一代入，四组偏振分别给
\begin{align*}
 K_{st}^{(x,\parallel)}
 &=\frac{4(1+c_\theta)}{1-c_\theta},&
 K_{st}^{(x,y)}
 &=2(1+c_\theta),\\
 K_{st}^{(y,\parallel)}
 &=2(1+c_\theta),&
 K_{st}^{(y,y)}
 &=\frac{4(1+c_\theta)}{1-c_\theta}.
\end{align*}
这里前、后两项中的 $1/(1-c_\theta)$ 来自 $t=-2Q^2(1-c_\theta)$；分子中的 $Q^2$ 则由两条外动量的 Dirac 迹提供，所以能量尺度消掉，只留下无量纲角函数。四项相加便得到
\begin{align*}
 K_{st}
 &=2(1+c_\theta)
 +2(1+c_\theta)
 +\frac{8(1+c_\theta)}{1-c_\theta}\\
 &=\frac{4(1+c_\theta)(3-c_\theta)}{1-c_\theta}.
\end{align*}
夸克传播子图与三胶子交换图的干涉因此归结为普通 Dirac 迹。物理干涉项还要乘颜色 Gram 元素 $G_{st}<0$，所以 $K_{st}$ 的符号本身不能判断最终干涉是增强还是相消。

剩余四个独立收缩使用同一套迹恒等式得到
\begin{align*}
 K_{uu}&=\frac{8(1+c_\theta^2)}{1+c_\theta},&
 K_{tt}&=\frac{8(1+c_\theta)(c_\theta^2-2c_\theta+5)}{(1-c_\theta)^2},\\
 K_{su}&=4(1-c_\theta),&
 K_{st}&=\frac{4(1+c_\theta)(3-c_\theta)}{1-c_\theta},&
 K_{ut}&=\frac{8(c_\theta^2-c_\theta+2)}{1-c_\theta}.
\end{align*}
由于 $K_{yx}=K_{xy}^*$，这六项已经覆盖三张图平方后的全部九个平方与干涉项。单个 $K_{xy}$ 依赖所选物理偏振基；与颜色结构相干组合并对完整物理振幅求和后，最终结果不依赖这一基底选择。

Lorentz--Dirac 收缩确定后，再独立计算颜色 Gram 矩阵。
把三张图的颜色矩阵记为
\begin{equation*}
 \mathcal C_s=T^bT^a,
 \qquad
 \mathcal C_u=T^aT^b,
 \qquad
 \mathcal C_t=\ii f^{abc}T^c=[T^a,T^b].
\end{equation*}
注意 $\mathcal C_t^\dagger=-\mathcal C_t$，因此干涉项的符号不能靠“颜色因子平方”猜测。定义
\begin{equation*}
 G_{xy}\equiv\sum_{a,b}\Tr_F(\mathcal C_x\mathcal C_y^\dagger).
\end{equation*}
由
\begin{equation*}
 \sum_aT^aT^a=C_FI,
 \qquad
 T^aT^bT^a=\left(C_F-\frac{C_A}{2}\right)T^b,
 \qquad
 \Tr(T^aT^b)=T_F\delta^{ab},
\end{equation*}
逐项得到
\begin{align*}
 G_{ss}=G_{uu}&=T_Fd_AC_F,\\
 G_{su}=G_{us}&=T_Fd_A\left(C_F-\frac{C_A}{2}\right),\\
 G_{tt}&=T_Fd_AC_A,\\
 G_{st}=G_{ts}&=-\frac12T_Fd_AC_A,\\
 G_{ut}=G_{tu}&=+\frac12T_Fd_AC_A,
\end{align*}
其中 $d_A=N_c^2-1$、$T_F=1/2$。例如
\begin{align*}
 G_{st}
 &=\sum_{a,b}\Tr\!\left[T^bT^a(T^bT^a-T^aT^b)\right]\\
 &=T_Fd_A\left(C_F-\frac{C_A}{2}\right)-T_Fd_AC_F
 =-\frac12T_Fd_AC_A,
\end{align*}
负的 $G_{st}$ 与正的 $G_{ut}$ 直接显示非阿贝尔交换图与两种夸克排序具有不同的颜色干涉符号。

最后先在振幅层相干求和，再实施初态平均。
未平均的完整平方振幅为
\begin{align*}
 \sum_{\rm spin,pol,color}|\mathcal M|^2
 =g_s^4\Bigl[&
 G_{ss}K_{ss}+G_{uu}K_{uu}+G_{tt}K_{tt}
 +2G_{su}K_{su}+2G_{st}K_{st}+2G_{ut}K_{ut}
 \Bigr].
\end{align*}
把上面六个 $K_{xy}$ 与六个颜色系数代入，并用
\begin{equation*}
 C_F=\frac{N_c^2-1}{2N_c},
 \qquad C_A=N_c,
 \qquad d_A=N_c^2-1,
\end{equation*}
整理成
\begin{equation*}
 \sum_{\rm spin,pol,color}|\mathcal M|^2
 =g_s^4\frac{d_A}{N_c}
 \frac{(c_\theta^2+2c_\theta+5)
 \left[N_c^2(c_\theta^2+2c_\theta+5)-(1-c_\theta)^2\right]}
 {(1-c_\theta)^2(1+c_\theta)}.
\end{equation*}
初态有 $2N_c$ 个夸克自旋--颜色态和 $2d_A$ 个胶子物理偏振--颜色态，所以平均因子是 $1/(4N_cd_A)$。再使用
\begin{align*}
 s^2+u^2&=4Q^4(c_\theta^2+2c_\theta+5),\\
 t^2&=4Q^4(1-c_\theta)^2,\\
 su&=-8Q^4(1+c_\theta),
\end{align*}
就得到
\begin{equation*}
 \overline{|\mathcal M(qg\to qg)|^2}
 =g_s^4(s^2+u^2)
 \left[
 \frac1{t^2}-\frac{C_F}{C_A}\frac1{su}
 \right],
\end{equation*}
即正文\eqnrefs{eq:qcd-qg-M2-dp94}。这里 $1/t^2$ 与 $1/(su)$ 是六类运动学收缩和六类颜色 Gram 元素相干组合后的结果；它们不能分别归属于某一张 Feynman 图的“概率”。
\end{derivation}

\begin{derivation}[从\eqnrefs{eq:qcd-gg-color-amplitude-dp94}到\eqnrefs{eq:qcd-gg-M2-dp94}]
这一过程若直接用 $-\eta_{\mu\nu}$ 展开四条外腿的偏振求和，会产生很长的 Lorentz 指标收缩，而且容易把“偏振代数”和“颜色代数”混在一起。这里改用物理螺旋度基底：先把每个非零螺旋度振幅算成三个颜色张量的线性组合，再单独完成一般 $SU(N_c)$ 的颜色求和。这样每一步都只处理一种数学结构。

先固定质心系运动学与外胶子的物理偏振。
令 $Q\equiv\sqrt{s}/2$，取度规 $\eta_{\mu\nu}=\operatorname{diag}(1,-1,-1,-1)$，并写
\begin{align*}
 k_1^\mu&=Q(1,0,0,1),&
 k_2^\mu&=Q(1,0,0,-1),\\
 k_3^\mu&=Q(1,\sin\theta,0,\cos\theta),&
 k_4^\mu&=Q(1,-\sin\theta,0,-\cos\theta).
\end{align*}
于是
\begin{equation*}
 t=-\frac{s}{2}(1-\cos\theta),
 \qquad
 u=-\frac{s}{2}(1+\cos\theta),
 \qquad s+t+u=0.
\end{equation*}
四个外胶子都只取横向物理偏振。以右螺旋度为例，可选
\begin{align*}
 \varepsilon_R^\mu(k_1)&=\frac1{\sqrt2}(0,1,\ii,0),&
 \varepsilon_R^\mu(k_2)&=\frac1{\sqrt2}(0,-1,\ii,0),\\
 \varepsilon_R^{*\mu}(k_3)&=\frac1{\sqrt2}(0,-\cos\theta,\ii,\sin\theta),&
 \varepsilon_R^{*\mu}(k_4)&=\frac1{\sqrt2}(0,\cos\theta,\ii,-\sin\theta).
\end{align*}
左螺旋度由上式作复共轭得到。逐一检查可得
\begin{equation*}
 k_i\cdot\varepsilon_{R,L}(k_i)=0,
 \qquad
 \varepsilon_{R,L}^*(k_i)\cdot\varepsilon_{R,L}(k_i)=-1,
\end{equation*}
所以这里从一开始就没有引入纵向或时间样偏振。

为缩短颜色指标，定义
\begin{equation*}
 \mathcal C_s\equiv f^{abe}f^{cde},
 \qquad
 \mathcal C_t\equiv f^{ace}f^{bde},
 \qquad
 \mathcal C_u\equiv f^{ade}f^{bce}.
\end{equation*}
Jacobi 恒等式给出
\begin{equation*}
 \mathcal C_s-\mathcal C_t+\mathcal C_u=0.
\end{equation*}

随后逐图计算独立的螺旋度振幅。
先算 $g_Rg_R\to g_Rg_R$。为避免与整体系数的无关相位纠缠，下表列出费曼规则给出的 $\ii\mathcal M$。三行交换图分别含两个三胶子顶角和一个 $1/s$、$1/t$、$1/u$ 传播子，最后一行为四胶子接触顶角。把上面的四个偏振矢量逐分量代入三胶子顶角
\begin{equation*}
 \Gamma^{\mu\nu\rho}(p,q,r)
 =\eta^{\mu\nu}(p-q)^\rho
 +\eta^{\nu\rho}(q-r)^\mu
 +\eta^{\rho\mu}(r-p)^\nu,
 \qquad p+q+r=0,
\end{equation*}
即可得到
\begin{center}
\begin{tabular}{c|c}
图 & $\ii\mathcal M(g_Rg_R\to g_Rg_R)/(\ii g_s^2)$\\ \hline
$s$ 道 & $-\mathcal C_s\cos\theta$\\[2pt]
$t$ 道 & $\displaystyle
 \mathcal C_t\frac{19+7\cos\theta-11\cos^2\theta+\cos^3\theta}
 {4(1-\cos\theta)}$\\[8pt]
$u$ 道 & $\displaystyle
 \mathcal C_u\frac{19-7\cos\theta-11\cos^2\theta-\cos^3\theta}
 {4(1+\cos\theta)}$\\[8pt]
四胶子接触项 & $\displaystyle
 -\mathcal C_s\cos\theta
 -\frac{\mathcal C_t}{4}(3+2\cos\theta-\cos^2\theta)
 -\frac{\mathcal C_u}{4}(3-2\cos\theta-\cos^2\theta)$
\end{tabular}
\end{center}
这里每一行只是有限个四维分量的乘加。以 $s$ 道为例，传播子分母为 $(k_1+k_2)^2=s$；两个三胶子顶角分别与
$\varepsilon_R(k_1)\varepsilon_R(k_2)$ 和
$\varepsilon_R^*(k_3)\varepsilon_R^*(k_4)$ 收缩。利用
$k_i\cdot\varepsilon_i=0$ 后，分子中的 $Q^2=s/4$ 恰好抵消传播子的 $1/s$，留下表中的 $-\mathcal C_s\cos\theta$。$t$、$u$ 道保留的 $1/(1\mp\cos\theta)$ 正是后来前向、后向交换极点的来源；接触项没有传播子，因此只给角度多项式。

把四行相加，并用 $\mathcal C_s=\mathcal C_t-\mathcal C_u$，分子多项式发生大规模相消：
\begin{align*}
 \frac{\ii\mathcal M_{RR\to RR}}{\ii g_s^2}
 &=4\left(
 \frac{\mathcal C_t}{1-\cos\theta}
 +\frac{\mathcal C_u}{1+\cos\theta}
 \right)\\
 &=-2\left(
 \mathcal C_t\frac{s}{t}
 +\mathcal C_u\frac{s}{u}
 \right).
\end{align*}
这一步显示了四胶子接触项的作用：它本身没有交换极点，却与三个交换图的多项式部分相消，使完整振幅只剩满足 Ward 恒等式与交叉结构的组合。

其余非零螺旋度振幅由同一分量计算或交叉变换得到。取三个独立通道可写成
\begin{align*}
 \mathcal M_{RR\to RR}
 &=-2g_s^2s\left(\frac{\mathcal C_t}{t}+\frac{\mathcal C_u}{u}\right),\\
 \mathcal M_{LR\to LR}
 &=2g_s^2u\left(\frac{\mathcal C_t}{t}+\frac{\mathcal C_s}{s}\right),\\
 \mathcal M_{LR\to RL}
 &=-2g_s^2t\left(\frac{\mathcal C_s}{s}-\frac{\mathcal C_u}{u}\right).
\end{align*}
整体相位不影响平方振幅。宇称变换给出另外三个等模振幅；其余螺旋度组合在树级为零。因此螺旋度求和只需把这三个独立结果平方后乘以 $2$；因子 $2$ 来自宇称相关的另外三个等模通道。

螺旋度核确定后，计算一般 $SU(N_c)$ 的颜色 Gram 矩阵。
记伴随表示维数
\begin{equation*}
 d_A=N_c^2-1.
\end{equation*}
首先，三个颜色张量的平方范数相同。例如
\begin{align*}
 \sum_{a,b,c,d}\mathcal C_s^2
 &=\sum_{e,e'}
 \left(\sum_{a,b}f^{abe}f^{ab e'}\right)
 \left(\sum_{c,d}f^{cde}f^{cd e'}\right)\\
 &=C_A^2\sum_{e,e'}\delta^{ee'}\delta^{ee'}
 =C_A^2d_A.
\end{align*}
令伴随表示生成元为 $(F^a)_{bc}=-\ii f^{abc}$。由
\begin{equation*}
 \Tr(F^aF^b)=C_A\delta^{ab},
 \qquad
 \sum_aF^aF^bF^a=\frac{C_A}{2}F^b,
\end{equation*}
可得一个代表性的交叉颜色积
\begin{align*}
 \sum_{a,b,c,d}\mathcal C_s\mathcal C_t
 &=\sum_{e,e'}\Tr(F^eF^{e'}F^eF^{e'})\\
 &=\frac{C_A}{2}\sum_{e'}\Tr(F^{e'}F^{e'})
 =\frac12C_A^2d_A.
\end{align*}
再结合 Jacobi 关系 $\mathcal C_s-\mathcal C_t+\mathcal C_u=0$ 及外颜色指标的置换对称性，得到完整 Gram 矩阵
\begin{equation*}
 G_{xy}\equiv\sum_{a,b,c,d}\mathcal C_x\mathcal C_y
 =C_A^2d_A
 \begin{pmatrix}
 1&\tfrac12&-\tfrac12\\
 \tfrac12&1&\tfrac12\\
 -\tfrac12&\tfrac12&1
 \end{pmatrix}_{xy},
 \qquad x,y\in\{s,t,u\}.
\end{equation*}
负的 $G_{su}$ 并不表示负概率；它只记录两个颜色张量在颜色向量空间中的夹角，最终概率由完整振幅的 Gram 二次型给出。

最后完成颜色求和、螺旋度求和与初态平均。
利用上面的 Gram 矩阵，第一个独立螺旋度通道给出
\begin{align*}
 \sum_{\rm color}|\mathcal M_{RR\to RR}|^2
 &=4g_s^4C_A^2d_A\,s^2
 \left(
 \frac1{t^2}+\frac1{u^2}+\frac1{tu}
 \right).
\end{align*}
同理
\begin{align*}
 \sum_{\rm color}|\mathcal M_{LR\to LR}|^2
 &=4g_s^4C_A^2d_A\,u^2
 \left(
 \frac1{t^2}+\frac1{s^2}+\frac1{ts}
 \right),\\
 \sum_{\rm color}|\mathcal M_{LR\to RL}|^2
 &=4g_s^4C_A^2d_A\,t^2
 \left(
 \frac1{s^2}+\frac1{u^2}+\frac1{su}
 \right).
\end{align*}
宇称伙伴再给相同贡献，因此全螺旋度、全颜色求和为
\begin{equation*}
 \sum_{\rm hel,color}|\mathcal M|^2
 =8g_s^4C_A^2d_A\,\mathcal S(s,t,u),
\end{equation*}
其中
\begin{align*}
 \mathcal S(s,t,u)={}&
 s^2\left(\frac1{t^2}+\frac1{u^2}+\frac1{tu}\right)
 +t^2\left(\frac1{s^2}+\frac1{u^2}+\frac1{su}\right)\\
 &+u^2\left(\frac1{s^2}+\frac1{t^2}+\frac1{st}\right).
\end{align*}
利用 $s+t+u=0$，例如
\begin{equation*}
 -\frac{2su}{t^2}
 =\frac{s^2+u^2}{t^2}-1,
\end{equation*}
以及循环置换，可把它化成
\begin{equation*}
 \mathcal S(s,t,u)
 =2\left(
 3-\frac{su}{t^2}-\frac{st}{u^2}-\frac{tu}{s^2}
 \right).
\end{equation*}
每个初态胶子有 $d_A$ 个颜色和两个物理螺旋度，所以初态简并度为 $4d_A^2$。最终
\begin{align*}
 \overline{|\mathcal M|^2}
 &=\frac{1}{4d_A^2}
 \sum_{\rm hel,color}|\mathcal M|^2\\
 &=\frac{4g_s^4C_A^2}{d_A}
 \left(
 3-\frac{su}{t^2}-\frac{st}{u^2}-\frac{tu}{s^2}
 \right),
\end{align*}
再代入 $d_A=N_c^2-1$，即得到正文\eqnrefs{eq:qcd-gg-M2-dp94}。
\end{derivation}
